\scanpage{8}

Supposons donc que $F_0$ satisfasse $\alpha)$ $\beta)$ $\gamma)$, \deleted{et soit $A\in F$, on a $A=\bigcup_i C_i$ (finie) $C_i=D_i-E_i$, $E_i$ non dense}

Alors $\deleted{\forall A,B\in F_0}\inserted{\forall A\in F}$, $\deleted{\overline{A\cap\complement B}}\inserted{\overline A}$ est le plus petit \uuline{$\in F_0$} contenant $A\cap\complement B$ (donc est connu en termes de $F_0$ d'entrée, sans appel à la top. de $X$ !). Donc la \deleted{\unsure{Cond}} condition pour $A\subset B$, $A,B\in F$, que $A$ non dense $B$ i.e. $A'=B-A$ dense en $B$ i.e. $\deleted{\unsure{B}}B\subset\overline{A'}$, a un sens indépendamment de la connaissance de la top. de $X$. L'axiome e) sur $F$ peut donc s'expliciter en termes de $F_0$, \unsure{disons} :

$\delta)$ $\forall$ suite décroissante $A_1\supset A_2\supset\ldots\supset A_n\supset\ldots$ d'élts de $F$, tels que $\forall n$, $A_{n+1}$ \deleted{non} ``$F_0$-\unsure{rare}'' dans $A_n$ (i.e. \deleted{$A_n$ \unsure{soit} $\varnothing$} tout $\in F_0$ qui contient $A_n-A_{n+1}$ contient $A_n$), on a $A_n=\varnothing$ pour $n$ assez grand.

\uline{NB} On peut se demander, \unsure{indépendt} d'une topologie de $X$, pour un ensemble $F_0$ de parties de $X$, quand il existe une top. sur $X$ pour laquelle $F_0$ \unsure{satisfais|satisfait} : $\alpha)$ $\beta)$ $\gamma)$ (et éventuellement $\delta)$). Il faut et s. que $F_0$ satisfasse $\alpha)$ $\beta)$ (qui ne font pas appel à une topologie de $X$), et :

\mspage{9}
